\documentclass[10pt,a4paper]{amsart}
\usepackage[margin=19mm]{geometry}
\usepackage{amsmath,amssymb,mathtools}
\usepackage[scr=rsfs,cal=euler]{mathalfa}
\usepackage{xfrac}
\usepackage[unicode,hidelinks,bookmarks=false]{hyperref}
\newcommand{\ie}{i.e.\ }
\newcommand{\cf}{cf.\ }
\newcommand{\resp}{respectively\ }
\newcommand{\Subsection}{\subsection}
\providecommand{\nobreakdash}{\nobreak\hbox{-}}
\setlength{\emergencystretch}{2em}
\multlinegap0pt
\usepackage{amssymb}
\usepackage[all]{xy}
\usepackage{xcolor}
\newtheorem{lemma}{Lemma}
\title{Boundedness of Fourier Integral Operators with complex phases on Fourier Lebesgue spaces}
\author{Duv\'an Cardona, William Obeng-Denteh, Frederick Opoku}
\date{Source: arXiv:2512.24854v2 (CC BY 4.0)}
\begin{document}
\maketitle

\noindent\emph{Source and licence.} Boundedness of Fourier Integral Operators with complex phases on Fourier Lebesgue spaces. Version arXiv:2512.24854v2 (CC BY 4.0), distributed by arXiv under the Creative Commons Attribution 4.0 International licence, http://creativecommons.org/licenses/by/4.0/, as recorded in the arXiv licence metadata for this exact version and shown on the abstract page https://arxiv.org/abs/2512.24854v2. Full licence notice: \texttt{licenses/arxiv-2512.24854-CC-BY-4.0.txt}.\par\smallskip

\noindent\textit{Editorial scope.} Counted unit: the article's lemma lemma (source0.tex lines 569--579). The article's complete original proof of it is delivered with it (source0.tex lines 580--644, one \texttt{proof} environment of 65 lines). No mathematics is written, repaired or re-derived: the statement and the proof are the article's own lines, in the article's own notation, and no retained line is substituted or re-flowed.  No further source-local context is needed: every cross-reference of every retained line resolves inside the two delivered spans.  Nothing was omitted from any retained span, so the package carries no omission record.  No retained line contains a citation, so the wrapper carries no bibliography and no bibliography entry.  No retained line contains a figure environment, an image-inclusion command or drawing code, so no image is delivered and the package declares no asset dependency.  Editorial formatting only, in addition: an AMS \texttt{amsart} preamble that restates the article's own declarations and macros used by the retained lines, namely the environments \texttt{lemma} on separate counters, together with the standard packages those lines need.  The retained environments print in this document as Lemma 1 in order of appearance, whereas the article prints the same items with the numbering of its own full document.  The complete native source of the article as distributed in the arXiv e-print for this exact version is bundled unchanged as \texttt{sources/191935/source0.tex}.
\begin{lemma}
    Let $\Phi_\tau(\eta,y)$ be a smooth phase function and, for fixed $j$ and $v$, consider the splitting $\eta = (\eta^{\prime},\eta^{\prime\prime})$, where $\eta^{\prime\prime}$ are tangent to a small cap $v$. Define the operator
\begin{align*}
L^v_j=(1-\langle2^{-j/2}\nabla_{{\eta}^{\prime}},2^{-j/2}\nabla_{{\eta}^{\prime}}\rangle)(1-\langle\nabla_{{\eta}^{\prime\prime}},\nabla_{{\eta}^{\prime\prime}}\rangle).
\end{align*}
Then $L_j^v$ satisfies the identity
\begin{align*}
     (1+4\pi^22^{-j}|( \nabla_1\Phi_{\tau}(\eta_j^v,y)-x)^{\prime}|^2)(1+4\pi^2|( \nabla_1\Phi_{\tau}(\eta_j^v,y)-x)^{\prime\prime}|^2)e^{2\pi i \langle \nabla_1\Phi_{\tau}(\eta_j^v,y)-x,\eta\rangle}\\
     =L^v_je^{2\pi i \langle \nabla_1\Phi_{\tau}(\eta_j^v,y)-x,\eta\rangle}.
\end{align*}
\end{lemma}

\begin{proof}
Let
\[
A:=\nabla_1\Phi_\tau(\eta_j^v,y)-x
\qquad\text{and}\qquad
e(\eta):=e^{2\pi i\langle A,\eta\rangle}
= e^{2\pi i\langle A',\eta'\rangle}\,e^{2\pi i\langle A'',\eta''\rangle},
\]
where $A=(A',A'')$ corresponds to the splitting of $\eta$. We compute the action of each factor in $L_j^v$ on the exponential.

\medskip\noindent\textbf{(i) Action on the $\eta'$--factor.}
For any coordinate $k$ in the $\eta'$ variables,
\[
\partial_{\eta'_k} e(\eta) = 2\pi i A_k' \, e(\eta),
\]
hence
\[
\partial_{\eta'_k}^2 e(\eta) = (2\pi i A_k')^2 e(\eta) = -4\pi^2 (A_k')^2 e(\eta).
\]
Therefore (using linearity and summing over $k$),
\[
\langle \nabla_{\eta'},\nabla_{\eta'}\rangle e(\eta)
= \sum_{k} \partial_{\eta'_k}^2 e(\eta)
= -4\pi^2 |A'|^2 e(\eta).
\]
Multiplying by the $2^{-j}$ factor from $L_j^v$ we obtain
\[
\langle 2^{-j/2}\nabla_{\eta'},2^{-j/2}\nabla_{\eta'}\rangle e(\eta)
= 2^{-j} \langle \nabla_{\eta'},\nabla_{\eta'}\rangle e(\eta)
= -4\pi^2 2^{-j} |A'|^2 e(\eta).
\]
Thus
\[
\big(1 - \langle 2^{-j/2}\nabla_{\eta'},2^{-j/2}\nabla_{\eta'}\rangle\big) e(\eta)
= \big(1 + 4\pi^2 2^{-j} |A'|^2\big) e(\eta).
\]

\medskip\noindent\textbf{(ii) Action on the $\eta''$--factor.}
Exactly the same computation (without the $2^{-j}$ weight) for the $\eta''$--variables gives
\[
\langle \nabla_{\eta''},\nabla_{\eta''}\rangle e(\eta)
= -4\pi^2 |A''|^2 e(\eta),
\]
and hence
\[
\big(1 - \langle \nabla_{\eta''},\nabla_{\eta''}\rangle\big) e(\eta)
= \big(1 + 4\pi^2 |A''|^2\big) e(\eta).
\]

\medskip\noindent\textbf{(iii) Combine the two factors.}
Since the $\eta'$ and $\eta''$ variables are independent and the two factors
in $L_j^v$ commute when acting on a product of functions that separate in
$\eta'$ and $\eta''$ (in particular on $e(\eta)$), we may apply them successively:
\[
L_j^v e(\eta)
= \big(1 - \langle 2^{-j/2}\nabla_{\eta'},2^{-j/2}\nabla_{\eta'}\rangle\big)
  \big(1 - \langle \nabla_{\eta''},\nabla_{\eta''}\rangle\big)
  e(\eta).
\]
By the two identities above we obtain,
\[ L_j^v e(\eta)=
\big(1 + 4\pi^2 2^{-j} |A'|^2\big)\big(1 + 4\pi^2 |A''|^2\big) e(\eta),
\]
which is exactly the claimed identity after recalling the definition of $A$.
\end{proof}

\end{document}
